\documentclass[10pt,a4paper]{amsart}
\usepackage[margin=19mm]{geometry}
\usepackage{amsmath,amssymb,mathtools}
\usepackage[scr=rsfs,cal=euler]{mathalfa}
\usepackage{xfrac}
\usepackage[unicode,hidelinks,bookmarks=false]{hyperref}
\newcommand{\ie}{i.e.\ }
\newcommand{\cf}{cf.\ }
\newcommand{\resp}{respectively\ }
\newcommand{\Subsection}{\subsection}
\providecommand{\nobreakdash}{\nobreak\hbox{-}}
\setlength{\emergencystretch}{2em}
\multlinegap0pt
\usepackage{xcolor}
\usepackage{amssymb}
\usepackage[normalem]{ulem}
\usepackage[shortlabels]{enumitem}
\usepackage{float}
\usepackage{dsfont}
\usepackage{csquotes}
\usepackage{algorithm}\usepackage{algpseudocode}
\usepackage{mathtools}
\newtheorem{lem}{Lemma}
\def\R{\mathbb{R}}
\title{Rapid decay and functional calculus in ${\rm C}^*$-probability spaces}
\author{Felipe Flores}
\date{Source: arXiv:2609.00452v1 (CC BY 4.0)}
\begin{document}
\maketitle

\noindent\emph{Source and licence.} Rapid decay and functional calculus in ${\rm C}^*$-probability spaces. Version arXiv:2609.00452v1 (CC BY 4.0), distributed by arXiv under the Creative Commons Attribution 4.0 International licence, http://creativecommons.org/licenses/by/4.0/, as recorded in the arXiv licence metadata for this exact version and shown on the abstract page https://arxiv.org/abs/2609.00452v1. Full licence notice: \texttt{licenses/arxiv-2609.00452-CC-BY-4.0.txt}.\par\smallskip

\noindent\textit{Editorial scope.} Counted unit: the article's lem iterated (source0.tex lines 506--517). The article's complete original proof of it is delivered with it (source0.tex lines 519--548, one \texttt{proof} environment of 30 lines). No mathematics is written, repaired or re-derived: the statement and the proof are the article's own lines, in the article's own notation, and no retained line is substituted or re-flowed.  No further source-local context is needed: every cross-reference of every retained line resolves inside the two delivered spans.  Nothing was omitted from any retained span, so the package carries no omission record.  No retained line contains a citation, so the wrapper carries no bibliography and no bibliography entry.  No retained line contains a figure environment, an image-inclusion command or drawing code, so no image is delivered and the package declares no asset dependency.  Editorial formatting only, in addition: an AMS \texttt{amsart} preamble that restates the article's own declarations and macros used by the retained lines, namely the environments \texttt{lem} on separate counters, together with the standard packages those lines need.  The retained environments print in this document as Lemma 1 in order of appearance, whereas the article prints the same items with the numbering of its own full document.  The complete native source of the article as distributed in the arXiv e-print for this exact version is bundled unchanged as \texttt{sources/209030/source0.tex}.
\begin{lem}\label{iterated}
Let $(A,\rho)$ be a tracial ${\rm C}^*$-probability space, and let
$L=(L_n)_{n=0}^\infty$ be a filtration of $A$.  If $a\in L_m$, then
\begin{equation*}
 p_jap_k=0
 \qquad\text{whenever}\qquad |j-k|>m.
\end{equation*}
Consequently, $\bigcup_{n\in\mathbb N}{L_n} \subset\bigcap_{k\in\mathbb N}{\rm dom}(\partial^k)$ and $ \|\partial^r(a)\|\leq m^r\|a\|$, when $a\in L_m$ and $r\geq1$. Furthermore, if $a\in A\cap{\rm dom}(\partial^r)$, then $a\in{\rm dom}((D_L-1)^r)$ and
\begin{equation*}
 \partial^r(a)\hat 1=i^r(D_L-1)^r\hat a=i^r\sum_{k\in\mathbb N}k^rp_k(\hat a).
\end{equation*}
\end{lem}

\begin{proof}
Let $a\in L_m$.  If $j>k+m$, then
$$
 a\bigl(p_kH\bigr)\subset a\overline{L_k}^{\,\|\cdot\|_2}\subset\overline{L_{m+k}}^{\,\|\cdot\|_2},
$$
and hence $p_jap_k=0$. On the other hand, if $k>j+m$, then $(p_jap_k)^*=p_ka^*p_j=0$, because $L_m$ is stable under $*$. This proves the first part of the lemma.

Now, for $-m\leq d\leq m$, set
\[
 a_d=\sum_{k\geq\max\{0,-d\}}p_{k+d}ap_k,
\]
where the sum converges in the strong operator topology. The previous observation implies that $a=\sum_{d=-m}^ma_d$ and
\begin{equation*}
 \beta_t(a)=\sum_{d=-m}^me^{idt}a_d,
\end{equation*}
which is a $\mathbb B(L^2(A,\rho))$-valued trigonometric polynomial of degree at most $m$. In particular, it is smooth as a function of $t$, which implies that $a\in \bigcap_{k\in\mathbb N}{\rm dom}(\partial^k) $. Furthermore, by the definition of the infinitesimal generator, we have that $\partial^r(a)=\frac{d^r}{dt^r}\beta_t(a)\big|_{t=0}$. Now, we apply the Bernstein inequality for trigonometric polynomials to see that
$$
\|\partial^r(a)\|\leq \sup_{t\in\R}\Big\|\frac{d^r}{dt^r}\beta_t(a) \Big\|\leq m^r \sup_{t\in\R}\|\beta_t(a)\|=m^r\|a\|.
$$

In order to prove the last assertion, we now suppose that $a\in A\cap{\rm dom}(\partial^r)$. Since $D_L(\hat 1)=\hat 1$, we have
\begin{equation*}
 \beta_t(a)\hat 1=U_taU_t^*\hat 1=e^{-it}U_t\hat a=e^{it(D_L-1)}\hat a.
\end{equation*}
Differentiating $r$ times, we get that
$$
 \partial^r(a)\hat1=\frac{d^r}{dt^r}\beta_t(a)\big|_{t=0}\hat 1=i^r(D_L-1)^r\hat a.
$$
Finally, $D_L-1$ acts as multiplication by $k$ on $p_kL^2(A,\rho)$, which gives the last identity.
\end{proof}

\end{document}
